\documentclass[11pt,a4paper]{article}

\usepackage{amsmath,amssymb,amsthm}
\usepackage{mathtools}
\usepackage{hyperref}
\usepackage{cleveref}
\usepackage[margin=1in]{geometry}
\usepackage{enumitem}
\usepackage{booktabs}

\newtheorem{theorem}{Theorem}[section]
\newtheorem{lemma}[theorem]{Lemma}
\newtheorem{proposition}[theorem]{Proposition}
\newtheorem{corollary}[theorem]{Corollary}
\newtheorem{definition}[theorem]{Definition}
\newtheorem{axiom}{Axiom}
\newtheorem{remark}[theorem]{Remark}
\newtheorem{example}[theorem]{Example}

\newcommand{\Prob}{\mathsf{Prob}}
\newcommand{\ProbProp}{\mathsf{ProbProp}}
\newcommand{\E}{\mathbb{E}}
\newcommand{\pand}{\wedge_p}
\newcommand{\por}{\vee_p}
\newcommand{\pnot}{\neg_p}
\newcommand{\pimplies}{\rightarrow_p}
\newcommand{\pzero}{\mathbb{0}}
\newcommand{\pone}{\mathbb{1}}
\newcommand{\ple}{\leq_p}
\newcommand{\padd}{+_p}
\newcommand{\pmul}{\cdot_p}
\newcommand{\psub}{-_p}

\title{Classical Logic as the $\{0,1\}$ Boundary Case of Probability\\[0.5em]
\large An Alternative Axiomatization of Mathematics}

\author{Probabilistic Foundations Project}

\date{\today}

\begin{document}

\maketitle

\begin{abstract}
We present an axiomatization of mathematics using probability as the primitive notion, with classical logic emerging as the degenerate $\{0,1\}$ boundary case. This is not a claim that probability is ``more fundamental'' than logic---both are equally valid starting points. Rather, we argue that the probabilistic axiomatization offers conceptual and technical advantages: uncertainty becomes the default rather than an afterthought, and certain constructions (developed in subsequent papers) become possible without requiring the Axiom of Choice. The classical embedding serves as a consistency check, confirming that standard mathematics is recoverable, but the interesting mathematics happens in the full $[0,1]$ setting.
\end{abstract}

\section{Introduction}

\subsection{Two Ways to Start}

Mathematics can be axiomatized starting from different primitives:

\begin{enumerate}[label=(\roman*)]
    \item \textbf{Logic-first}: Begin with classical logic ($\land, \lor, \neg$, LEM), build set theory, derive probability as measure theory.
    \item \textbf{Probability-first}: Begin with probability ($\Prob$, $\E[\cdot|\cdot]$), derive classical logic as the $\{0,1\}$ restriction.
\end{enumerate}

Neither approach is ``more fundamental.'' They are dual axiomatizations. However, the probability-first approach has practical advantages:

\begin{itemize}
    \item Uncertainty is native, not bolted on
    \item Graded truth (confidence levels) is primitive
    \item Certain existence proofs become constructive (Paper 4)
    \item The framework naturally supports probabilistic reasoning
\end{itemize}

\subsection{What This Paper Does and Does Not Claim}

\textbf{We claim}: Classical logic is isomorphic to the $\{0,1\}$-valued fragment of our probabilistic system. The logical operations $\land, \lor, \neg$ correspond exactly to probabilistic operations restricted to extreme values.

\textbf{We do not claim}: That probability is ``more primitive'' or ``more fundamental'' than logic in any absolute sense. Our axioms for $\Prob$ are just as axiomatic as the axioms of classical logic.

\textbf{The payoff}: The interesting work is not the classical embedding (which is merely a sanity check), but what happens in the full $[0,1]$ setting---probabilistic proof theory (Paper 2), probabilistic natural numbers (Paper 3), and choice-free existence proofs (Paper 4).

\subsection{Machine Verification}

All results in this paper are formalized in Lean 4 and machine-verified. The implementation is available at \url{https://github.com/[repository]}.

\section{The Probability Type}

\subsection{Axiomatization}

We axiomatize an abstract type $\Prob$ representing probability values in $[0,1]$.

\begin{axiom}[Prob Type]
There exists a type $\Prob$ with distinguished elements $\pzero, \pone : \Prob$ and operations:
\begin{align*}
    \padd &: \Prob \to \Prob \to \Prob \\
    \pmul &: \Prob \to \Prob \to \Prob \\
    \psub &: \Prob \to \Prob \to \Prob \\
    \ple &: \Prob \to \Prob \to \mathsf{Prop}
\end{align*}
\end{axiom}

\begin{axiom}[Ordered Semiring Structure]
The type $\Prob$ forms an ordered semiring:
\begin{enumerate}[label=(\alph*)]
    \item $(\Prob, \padd, \pzero)$ is a commutative monoid
    \item $(\Prob, \pmul, \pone)$ is a commutative monoid
    \item Multiplication distributes over addition
    \item $\pzero \pmul p = \pzero$ for all $p$
    \item $\ple$ is a total order with $\pzero \ple p \ple \pone$ for all $p$
    \item Operations are monotonic with respect to $\ple$
\end{enumerate}
\end{axiom}

\begin{axiom}[Complement Structure]
For subtraction (truncated at $\pzero$):
\begin{enumerate}[label=(\alph*)]
    \item $p \psub p = \pzero$
    \item $(p \padd q) \psub q = p$
    \item $p \padd (\pone \psub p) = \pone$
    \item $\pone \psub (\pone \psub p) = p$
\end{enumerate}
\end{axiom}

\begin{axiom}[Non-triviality]
$\pzero \neq \pone$.
\end{axiom}

\begin{remark}
These axioms do not construct $\Prob$; they characterize it abstractly. A synthetic construction from type-theoretic primitives is the subject of Paper 5.
\end{remark}

\subsection{Interpretation}

Elements of $\Prob$ represent degrees of belief, frequencies, or abstract probability values. The key insight is that we work with these values \emph{directly}, without first building measure theory on top of set theory on top of logic.

\section{Probabilistic Propositions}

\subsection{Definition}

\begin{definition}[ProbProp]
A \emph{probabilistic proposition} over a type $\alpha$ is a function:
\[
\ProbProp(\alpha) := \alpha \to \Prob
\]
\end{definition}

Unlike classical propositions (which map to $\{\mathsf{true}, \mathsf{false}\}$), probabilistic propositions assign a probability value to each state. This is \emph{not} the probability that the proposition is true---it \emph{is} the truth value, generalized from Boolean to $[0,1]$.

\begin{example}
Let $\alpha$ be the type of coin flip outcomes $\{\mathsf{heads}, \mathsf{tails}\}$. A fair coin defines:
\[
\mathsf{fair}(x) = 0.5 \quad \text{for all } x
\]
A biased coin (70\% heads) defines:
\[
\mathsf{biased}(x) = \begin{cases} 0.7 & x = \mathsf{heads} \\ 0.3 & x = \mathsf{tails} \end{cases}
\]
\end{example}

\section{Probabilistic Logical Operations}

\subsection{Definitions}

We define logical operations on $\ProbProp$ using probability algebra:

\begin{definition}[Probabilistic Connectives]
For $A, B : \ProbProp(\alpha)$:
\begin{align}
    (A \pand B)(x) &:= A(x) \pmul B(x) \label{eq:pand} \\
    (A \por B)(x) &:= (A(x) \padd B(x)) \psub (A(x) \pmul B(x)) \label{eq:por} \\
    (\pnot A)(x) &:= \pone \psub A(x) \label{eq:pnot} \\
    (A \pimplies B)(x) &:= \pone \psub (A(x) \pmul (\pone \psub B(x))) \label{eq:pimplies}
\end{align}
\end{definition}

\begin{remark}
Equation \eqref{eq:por} is the inclusion-exclusion principle: $P(A \cup B) = P(A) + P(B) - P(A \cap B)$. Equation \eqref{eq:pand} assumes independence; for dependent propositions, conditional expectation (Paper 2) is required.
\end{remark}

\subsection{Algebraic Properties}

\begin{theorem}[Commutativity]
For all $A, B : \ProbProp(\alpha)$ and $x : \alpha$:
\begin{align*}
    (A \pand B)(x) &= (B \pand A)(x) \\
    (A \por B)(x) &= (B \por A)(x)
\end{align*}
\end{theorem}

\begin{proof}
Follows from commutativity of $\pmul$ and $\padd$ in $\Prob$.
\end{proof}

\begin{theorem}[Associativity of Conjunction]
$(A \pand (B \pand C))(x) = ((A \pand B) \pand C)(x)$
\end{theorem}

\begin{proof}
Follows from associativity of $\pmul$.
\end{proof}

\begin{theorem}[Identity and Absorption]
Let $\mathbf{0}(x) = \pzero$ and $\mathbf{1}(x) = \pone$ for all $x$. Then:
\begin{align*}
    (A \pand \mathbf{1})(x) &= A(x) & (A \pand \mathbf{0})(x) &= \pzero \\
    (A \por \mathbf{0})(x) &= A(x) & (A \por \mathbf{1})(x) &= \pone
\end{align*}
\end{theorem}

\begin{theorem}[De Morgan Laws]
\begin{align*}
    \pnot(A \pand B) &= (\pnot A) \por (\pnot B) \\
    \pnot(A \por B) &= (\pnot A) \pand (\pnot B)
\end{align*}
\end{theorem}

\section{The Law of Excluded Middle as Algebra}

\subsection{LEM is an Algebraic Identity}

In classical logic, the Law of Excluded Middle ($\phi \lor \neg\phi$) is an axiom or a consequence of other axioms. In our framework, it is an \emph{algebraic identity}:

\begin{theorem}[Probabilistic LEM]
For any $\phi : \ProbProp(\alpha)$ and $x : \alpha$:
\[
\phi(x) \padd (\pnot\phi)(x) = \pone
\]
\end{theorem}

\begin{proof}
By definition of $\pnot$:
\[
\phi(x) \padd (\pone \psub \phi(x)) = \pone
\]
This is Axiom (Complement Structure, part c).
\end{proof}

\begin{remark}
This is not ``LEM holds probabilistically.'' It is: \emph{LEM is the statement that probabilities sum to 1}. The logical law is a shadow of probability algebra.
\end{remark}

\subsection{Double Negation}

\begin{theorem}[Double Negation Elimination]
$(\pnot(\pnot\phi))(x) = \phi(x)$
\end{theorem}

\begin{proof}
$\pone \psub (\pone \psub \phi(x)) = \phi(x)$ by Axiom (Complement Structure, part d).
\end{proof}

\section{Classical Logic as the $\{0,1\}$ Restriction}

\subsection{Classical Propositions}

\begin{definition}[Classical Value]
A probability $p : \Prob$ is \emph{classical} if $p = \pzero$ or $p = \pone$.
\end{definition}

\begin{definition}[Classical Proposition]
A probabilistic proposition $\phi : \ProbProp(\alpha)$ is \emph{classical} if $\phi(x)$ is classical for all $x : \alpha$.
\[
\mathsf{ClassicalProp}(\alpha) := \{ \phi : \ProbProp(\alpha) \mid \forall x,\, \phi(x) \in \{\pzero, \pone\} \}
\]
\end{definition}

\subsection{Closure Theorems}

The classical propositions form a Boolean subalgebra:

\begin{theorem}[Classical AND is Classical]
If $A, B : \mathsf{ClassicalProp}(\alpha)$, then $A \pand B : \mathsf{ClassicalProp}(\alpha)$.
\end{theorem}

\begin{proof}
For each $x$, we have four cases:
\begin{center}
\begin{tabular}{cc|c}
$A(x)$ & $B(x)$ & $(A \pand B)(x) = A(x) \pmul B(x)$ \\
\midrule
$\pzero$ & $\pzero$ & $\pzero \pmul \pzero = \pzero$ \\
$\pzero$ & $\pone$ & $\pzero \pmul \pone = \pzero$ \\
$\pone$ & $\pzero$ & $\pone \pmul \pzero = \pzero$ \\
$\pone$ & $\pone$ & $\pone \pmul \pone = \pone$
\end{tabular}
\end{center}
In all cases, the result is classical.
\end{proof}

\begin{theorem}[Classical OR is Classical]
If $A, B : \mathsf{ClassicalProp}(\alpha)$, then $A \por B : \mathsf{ClassicalProp}(\alpha)$.
\end{theorem}

\begin{proof}
Similar case analysis using the definition $(A \por B)(x) = (A(x) \padd B(x)) \psub (A(x) \pmul B(x))$.
\end{proof}

\begin{theorem}[Classical NOT is Classical]
If $A : \mathsf{ClassicalProp}(\alpha)$, then $\pnot A : \mathsf{ClassicalProp}(\alpha)$.
\end{theorem}

\begin{proof}
If $A(x) = \pzero$, then $(\pnot A)(x) = \pone \psub \pzero = \pone$.
If $A(x) = \pone$, then $(\pnot A)(x) = \pone \psub \pone = \pzero$.
\end{proof}

\begin{theorem}[Classical Idempotence]
For classical $A$: $(A \pand A)(x) = A(x)$.
\end{theorem}

\begin{proof}
If $A(x) = \pzero$: $\pzero \pmul \pzero = \pzero$.
If $A(x) = \pone$: $\pone \pmul \pone = \pone$.
\end{proof}

\begin{remark}
Idempotence fails for non-classical propositions: if $A(x) = 0.5$, then $(A \pand A)(x) = 0.25 \neq 0.5$.
\end{remark}

\subsection{The Embedding is a Consistency Check}

The closure theorems show that $\mathsf{ClassicalProp}(\alpha)$ with $(\pand, \por, \pnot)$ forms a Boolean algebra isomorphic to the classical propositional logic over $\alpha$.

\textbf{This is a sanity check, not the goal.} It confirms that our probabilistic axiomatization is consistent with classical mathematics---anything provable classically remains provable in the $\{0,1\}$ fragment.

\textbf{The interesting mathematics} happens when we leave the $\{0,1\}$ boundary and work in the full $[0,1]$ setting. Papers 2--4 develop this.

\section{The Meta/Object Distinction}

\subsection{Three Levels}

Our formalization maintains a strict separation:

\begin{enumerate}
    \item \textbf{Metalogic} (Lean's type theory): Used to \emph{define} and \emph{verify} the probabilistic system. Lean's $\land, \lor, \neg, \forall, \exists$ appear here.

    \item \textbf{Object system} (Probabilistic foundations): Where actual mathematics happens. Only $\Prob$, $\E[\cdot|\cdot]$, $\pand$, $\por$, $\pnot$ appear here.

    \item \textbf{Classical fragment}: The $\{0,1\}$ restriction of the object system. Isomorphic to classical logic.
\end{enumerate}

\subsection{Avoiding Logic Slippage}

A common error is to use metalogic ($\forall, \exists$) at the object level. For example:

\textbf{Wrong}: ``$\forall x, P(x) \to Q(x)$'' as an object-level statement.

\textbf{Right}: ``$P \pimplies Q$'' as an object-level statement, or ``$\E[Q|P] = \pone$'' for entailment.

The metalogic quantifiers are for \emph{defining} the object system, not for reasoning \emph{within} it.

\section{Implementation}

\subsection{Lean 4 Code}

The core definitions in Lean:

\begin{verbatim}
axiom Prob : Type
axiom Prob.zero : Prob
axiom Prob.one : Prob
axiom Prob.le : Prob -> Prob -> Prop
axiom Prob.add : Prob -> Prob -> Prob
axiom Prob.mul : Prob -> Prob -> Prob
axiom Prob.sub : Prob -> Prob -> Prob

notation "0" => Prob.zero
notation "1" => Prob.one

def ProbProp (a : Type) := a -> Prob

noncomputable def prob_and (A B : ProbProp a) : ProbProp a :=
  fun x => A x * B x

noncomputable def prob_or (A B : ProbProp a) : ProbProp a :=
  fun x => (A x + B x) - (A x * B x)

noncomputable def prob_not (A : ProbProp a) : ProbProp a :=
  fun x => 1 - A x

def isClassical (p : Prob) : Prop := p = 0 \/ p = 1

def ClassicalProp (a : Type) :=
  { f : ProbProp a // forall x, isClassical (f x) }
\end{verbatim}

\subsection{Verification Status}

\begin{itemize}
    \item All theorems in this paper are machine-verified
    \item Build: \texttt{cd lean \&\& lake build}
    \item Zero warnings, zero errors
\end{itemize}

\section{Discussion}

\subsection{Philosophical Status}

We have presented probability as an \emph{alternative} primitive, not a \emph{superior} one. The advantages of this axiomatization are practical:

\begin{itemize}
    \item Graded truth is built-in, not retrofitted
    \item Probabilistic reasoning is native
    \item Certain constructions become possible without AC (Paper 4)
\end{itemize}

\subsection{Relation to Bayesian Probability}

Our framework is compatible with both frequentist and Bayesian interpretations. The $\Prob$ type is abstract---it could represent:
\begin{itemize}
    \item Long-run frequencies
    \item Degrees of belief
    \item Abstract weights in a categorical setting
\end{itemize}

\subsection{What Comes Next}

\begin{itemize}
    \item \textbf{Paper 2}: Probabilistic proof theory with $\Gamma \vdash_p \phi$ (``derive $\phi$ from $\Gamma$ with confidence $p$'')
    \item \textbf{Paper 3}: Natural numbers as distributions over $\mathbb{N}$
    \item \textbf{Paper 4}: Zorn's lemma via distributions, without the Axiom of Choice
    \item \textbf{Paper 5}: Synthetic construction of $\Prob$ from type-theoretic primitives
\end{itemize}

\section{Conclusion}

Classical logic is the $\{0,1\}$ boundary case of probability algebra. This observation, while mathematically straightforward, suggests a different way to organize foundational mathematics---one where uncertainty is primitive and certainty is the special case.

The classical embedding is a consistency check. The real work begins when we leave the boundary and develop mathematics in the full probabilistic setting.

\appendix

\section{Complete Axiom List}

For reference, the complete list of axioms for $\Prob$:

\begin{enumerate}
    \item $\pzero \ple p$ for all $p$
    \item $p \ple \pone$ for all $p$
    \item $p \ple p$ (reflexivity)
    \item $p \ple q \land q \ple r \Rightarrow p \ple r$ (transitivity)
    \item $p \ple q \land q \ple p \Rightarrow p = q$ (antisymmetry)
    \item $p \pmul q \ple p' \pmul q'$ when $p \ple p'$ and $q \ple q'$
    \item $\pone \pmul p = p$
    \item $p \pmul \pone = p$
    \item $p \pmul q = q \pmul p$
    \item $(p \pmul q) \pmul r = p \pmul (q \pmul r)$
    \item $\pzero \pmul p = \pzero$
    \item $p \pmul \pzero = \pzero$
    \item $p \padd q = q \padd p$
    \item $(p \padd q) \padd r = p \padd (q \padd r)$
    \item $\pzero \padd p = p$
    \item $p \padd \pzero = p$
    \item $p \psub p = \pzero$
    \item $(p \padd q) \psub q = p$
    \item $\pone \psub \pzero = \pone$
    \item $\pone \psub \pone = \pzero$
    \item $p \psub \pzero = p$
    \item $p \padd (\pone \psub p) = \pone$
    \item $\pone \psub (\pone \psub p) = p$
    \item $p \pmul (q \padd r) = (p \pmul q) \padd (p \pmul r)$
    \item $(p \padd q) \pmul r = (p \pmul r) \padd (q \pmul r)$
    \item $(\pone \padd \pone) \psub \pone = \pone$
    \item $\pzero \neq \pone$
\end{enumerate}

\bibliographystyle{plain}

\end{document}
